\subsection{Beskrivende dataanalyse (5 poeng)}\label{beskrivende-dataanalyse-5-poeng}

Her bør du gjennomføre en kort deskriptiv dataanalyse for å få en
oversikt over datasettet. Det betyr at man lager noen visualiseringer
for å forstå dataene bedre, og at man beregner noen enkle ting. Fokuser
på det som er relevant for modelleringen, og bruk sunn fornuft. For mye
detaljer eller figurer uten forklaring og tydelige akseetiketter er ikke
nyttige for leseren.

Du kan for eksempel:

\begin{itemize}
\tightlist
\item
  Presentere deskriptiv statistikk, som gjennomsnitt, median,
  standardavvik og varians, for å beskrive datasettets sentraltendens,
  variasjon og fordeling.
\item
  Bruke spredningsdiagrammer (scatter plots) og korrelasjonsmatriser for
  å undersøke sammenhenger mellom ulike variabler.
\item
  Bruke histogrammer for å vise fordelingen til relevante variabler.
\item
  Bruke boksplott (box plots) for å undersøke fordeling, variasjon og
  eventuelle uteliggere.
\end{itemize}

Resultatene bør forklares i teksten, med særlig fokus på funn som kan
være relevante for den videre modelleringen. Skriv en kort legende under
hver grafikk, slik at det blir tydelig hva grafikken viser.
